\documentclass[12pt]{article}
\usepackage[margin=1in]{geometry}
\usepackage{amsmath,amsthm,amssymb,scrextend}
\setlength{\parindent}{0pt}
\begin{document}
\newcommand{\ds}{\displaystyle}


\textbf{(i)} Consider the first two tosses

If the first toss is $\mathrm{T}$, two consecutive heads must need for either to win, which $THH$ occurs first

If they are $\mathrm{HH}$, then the first later tail gives $\mathrm{HHT}$, so $A$ wins

If they are $\mathrm{HT}$, similarly to the first reason, $B$ will win

The first two tosses are equally likely to be $\mathrm{HH},\mathrm{HT},\mathrm{TH},\mathrm{TT}$

$$P(A\text{ wins})=\frac14(1+0+0+0)=\boxed{\frac14}$$

\bigskip\hrule\bigskip

\textbf{(ii)} Consider the first two tosses

If they are $\mathrm{HH}$, then the first later tail gives $\mathrm{HHT}$, so $A$ will win

If they are $\mathrm{TT}$, then the first later head gives $\mathrm{TTH}$, so $C$ will win

If they are $\mathrm{HT}$ or $\mathrm{TH}$, only $B$ or $D$ can win

By symmetry, the total probabilities of $B$ and $D$ winning from these two cases are equal

$$\therefore \boxed{P(A)=P(B)=P(C)=P(D)=\frac14}$$


\bigskip\hrule\bigskip

\textbf{(iii)} If the first two tosses are $\mathrm{TT}$, then $\mathrm{T}$ loops and the first later $\mathrm{H}$ gives $\mathrm{TTH}$

$$\boxed{P(C\text{ wins}\mid \mathrm{TT})=1}$$

$$p=P(C\mid \mathrm{HT}),\qquad q=P(C\mid \mathrm{TH}),\qquad r=P(C\mid \mathrm{HH})$$

From $\mathrm{HT}$, the next toss is either $\mathrm{T}$, giving $\mathrm{TT}$, or $\mathrm{H}$, giving $\mathrm{TH}$

$$p=\frac12\cdot1+\frac12q$$

From $\mathrm{TH}$, the next toss is either $\mathrm{H}$, giving $\mathrm{THH}$, so $B$ wins, or $\mathrm{T}$, giving $\mathrm{HT}$

$$q=\frac12\cdot0+\frac12p=\frac12p$$

From $\mathrm{HH}$, the next toss is either $\mathrm{T}$, giving $\mathrm{HT}$, or $\mathrm{H}$, staying at $\mathrm{HH}$

$$r=\frac12p+\frac12r$$

$$\implies p=\frac23,\qquad q=\frac13,\qquad r=\frac23$$

The first two tosses are equally likely to be $\mathrm{HH},\mathrm{HT},\mathrm{TH},\mathrm{TT}$

\begin{align*}
P(C\text{ wins})
&=\frac14(r+p+q+1)\\
&=\frac14\left(\frac23+\frac23+\frac13+1\right)\\
&=\boxed{\frac23}
\end{align*}


\bigskip
\hrule
\begin{center}
    This solution was submitted by \textbf{Jon}\\[4pt]
    Published and refined by \textbf{Jon}\\[8pt]
    \textit{Last updated: June 24, 2026}
\end{center}

\end{document}
